\documentclass[11pt]{article}
\usepackage[margin=0.68in]{geometry}
\usepackage{booktabs,enumitem,fancyhdr,graphicx,overpic,tabularx,titlesec}
\usepackage[table]{xcolor}
\usepackage[hidelinks]{hyperref}
\graphicspath{{docs/lessons/assets/}{../assets/}}
\hypersetup{pdftitle={ADK Lesson 042: Tabletop Course Marshal},
 pdfauthor={Kevin Shanahan},
 pdfsubject={Explicit start authority, ordered checkpoints, and bounded volatile records},
 pdflang={en-US}}
\pagestyle{fancy}\fancyhf{}
\lhead{\textsc{ADK Lesson 042}}\rhead{Tabletop Course Marshal}
\cfoot{\thepage}\setlength{\headheight}{14pt}
\setlength{\parindent}{0pt}\setlength{\parskip}{5pt}
\setlist{nosep,leftmargin=1.45em}
\titleformat{\section}{\Large\bfseries}{\thesection}{0.5em}{}
\newcommand{\lessonbox}[1]{\begin{center}\fbox{\parbox{0.92\linewidth}{#1}}\end{center}}

\begin{document}
\begin{center}
 {\Huge\bfseries Tabletop Course Marshal}\\[4pt]
 {\Large Authorize once, cross in order, freeze the evidence}\\[8pt]
 \textit{E0 host verified; powered fixture and bench acceptance open}
\end{center}
\begin{tabularx}{\linewidth}{@{}>{\bfseries}l X >{\bfseries}l X@{}}
Lesson & 042 --- project & Time & 150--210 minutes\\
Prerequisites & Lessons 040--041 & Board & deterministic Mega replay\\
Evidence & checkpoint LEDs, heartbeat, all-red, display &
Boundary & hand-moved card; no actuator\\
\end{tabularx}

\lessonbox{\textbf{Authority rule.} PIR motion can make a start eligible, but
it can never start a run. Only a debounced press from the configured button,
in the same course-clock frame as fresh valid PIR Motion evidence, emits the
start event. A press outside that frame is not remembered.}

\section{A course made from semantic evidence}
\begin{center}
% ADK visual: pencil
\begin{overpic}[width=0.96\linewidth]{042-course-layout-pencil.png}
 \put(4,62){\small explicit button}
 \put(7,39){\small PIR eligibility}
 \put(26,61){\small checkpoint 1}
 \put(44,62){\small checkpoint 2}
 \put(61,56){\small checkpoint 3}
 \put(77,50){\small checkpoint 4}
 \put(87,38){\small finish}
 \put(15,5){\small 4 checkpoints + all-red}
 \put(64,5){\small bounded display}
\end{overpic}

\small\textit{Alternative text: a pencil drawing shows a finger pressing the
separate start button, a nearby PIR eligibility witness, a hand-moved card
crossing four ordered gates and a finish, then checkpoint LEDs, an all-red
indicator, and a display. The policy supports four checkpoints; the canonical
replay configures three.}
\end{center}

The project consumes qualified copied values rather than pins or retail
modules. Configuration maps one through four unique checkpoint IDs to optical
source kind, source ID, and calibration revision. That semantic binding, not
input-array or pin order, defines the course.

\texttt{CourseStartPolicy} receives the copied button edge and PIR state in
one frame. \texttt{CourseMarshal} receives its start event, one immutable
\texttt{PresenceSnapshot}, and at most four qualified checkpoint events. It
validates the entire frame before mutation. Changed same-time, backward, and
half-range-invalid frames fault without accepting a partial checkpoint.

\section{Walk the evidence timeline}
\begin{center}
% ADK visual: pencil
\begin{overpic}[width=0.98\linewidth]{042-run-timeline-pencil.png}
 \put(3,53){\small eligible}
 \put(23,53){\small authorize}
 \put(43,53){\small ordered checkpoints}
 \put(76,53){\small finish agreement}
 \put(89,48){\small frozen record}
 \put(24,3){\small simultaneous edges reject; array order cannot decide}
\end{overpic}

\small\textit{Alternative text: a pencil timeline separates PIR eligibility
from the required button press, continues through the configured checkpoints
and finish agreement to a record, and branches to all-red when gates collide.
The policy supports four checkpoints; the canonical replay configures three.}
\end{center}

Start in \texttt{Disarmed}. A structurally valid frame enters
\texttt{Arming}; fresh valid PIR Motion evidence makes the marshal
\texttt{Ready}. PIR alone leaves it there. Only a same-frame explicit start
event enters \texttt{Running}. A press before or after eligibility does
nothing and is never latched.

Each accepted checkpoint advances one semantic slot. Two distinct legal
events whose timestamps are inside the inclusive simultaneity window reject
as ambiguous. A repeated identity is idempotent, while a distinct event for an
already accepted checkpoint rejects as duplicate.

Finish acceptance requires the complete checkpoint prefix, a fresh valid
finish-guard activation event, and fresh valid approach-range evidence within
the agreement window. The resulting terminal record stays immutable until
\texttt{acknowledgeRecord()} or reconstruction.

\section{Fixed storage is an architectural choice}
The original monolithic design measured 653 bytes for the marshal and 432
bytes for its presenter on AVR, exceeding the 128-byte largest-object ceiling.
The remediated design uses six explicit caller-owned fixed storages:

\begin{tabularx}{\linewidth}{@{}X r@{}}
\toprule
Budget item & AVR bytes\\
\midrule
run, trigger, trigger-presence, and three replay storages & 456\\
six storages plus marshal & 525\\
storages plus marshal, start policy, and presenter & 639\\
\bottomrule
\end{tabularx}

Every individual public type and policy object now measures at most 128 bytes.
The marshal retains references to all six storages for its full lifetime;
input-view pointers are borrowed only for \texttt{update()}. There is no heap,
RTC, SD card, or persistence claim. Reconstruction begins a new volatile
identity epoch. Sequence exhaustion is permanent for one object lifetime.

These are E0 policy subtotals, not a complete Mega budget. The canonical
replay sketch measures 15,596 bytes of flash and 901 bytes of static SRAM on
the Mega 2560. Powered composition still requires aggregate current, pin,
cadence, and stack evidence.

\section{Observe without Serial}
\begin{tabularx}{\linewidth}{@{}l X@{}}
\toprule
Visible path & Prediction and interpretation\\
\midrule
three checkpoint LEDs & canonical replay evidence; policy supports four\\
ready/heartbeat & eligible readiness and supplied-time progress\\
all-red & rejected or faulted run; inspect the retained trigger\\
four-digit display & one bounded cell per update; elapsed or fault code\\
\bottomrule
\end{tabularx}

\texttt{CourseMarshalPresenter} copies only the compact marshal snapshot. It
advances at most one display cell per later timestamp, so a sparse update
cannot monopolize the loop. Display or renderer failure cannot alter course
phase, order, elapsed time, or the frozen record.

\section{Predict, observe, interpret}
\begin{enumerate}
 \item Supply valid PIR Motion evidence without a button press. Predict
       \texttt{Ready}, not \texttt{Running}.
 \item Press the configured button in the same eligible frame. Predict exactly
       one start event. Repeat a held input and predict no second event.
 \item Cross checkpoint IDs in configured order. Predict one additional local
       LED after each accepted event.
 \item Present two distinct checkpoint events within the simultaneity window.
       Permute their array positions; predict the same rejection and all-red.
 \item Reach the finish guard early or without fresh approach evidence.
       Predict rejection with trigger evidence preserved.
 \item Complete the ordered course and agreeing finish. Predict a stable
       accepted record and elapsed display until acknowledgment.
 \item Inject failed source status or malformed time. Predict a fault without
       partially advancing the accepted prefix.
\end{enumerate}

\section{Powered boundary and shutdown}
Lesson 042 publishes no wiring table and no electrically authoritative
schematic. Exact optical, PIR, HC-SR04, button, display, and indicator
specimens must first close identity, pin order, voltage, output-rail, polarity,
pull-up, emitter, current, resource, timing, stack, and aggregate SRAM gates.
Only then may a separately reviewed E1 composition add pin-by-pin prose and a
formal schematic.

The physical model is moved only by hand. There is no motor, powered gate,
vehicle control, launcher control, or stored-energy command. Disconnect USB
before wiring changes and stop for heat, odor, reset, unstable rail, or
unexpected output.

Resource acquisition and electrical safe state are separate evidence. A
future qualified fixture must prove claims and rollback independently from
inactive indicators and high-impedance sensor endpoints after shutdown.

\lessonbox{\textbf{Physical acceptance remains open.} Compilation, replay,
object-size measurements, pencil drawings, and documentation do not prove a
sensor, current budget, display, timing cadence, stack peak, or safe shutdown
on a bench. Leave the acceptance card blank until those observations exist.}

\section{Software evidence}
Host tests cover authorization timing, eligibility phases, checkpoint order,
simultaneity, replay identity, finish agreement, timeout, sequence exhaustion,
record acknowledgment, presenter pacing, rollover, reset, and source faults.
The canonical Mega sketch is a deterministic copied-evidence replay; it is not
a claim that any powered module is qualified.

\begin{itemize}
 \item ADK Lessons 040--042 optical course marshal design record.
 \item ADK deterministic testing contract and safety model.
 \item ADK component architecture stress-pass template.
\end{itemize}

This untagged print edition has a searchable HTML alternative.
It makes no PDF/UA or physical accessibility claim.
\end{document}
